
\documentclass[11pt]{article}
\usepackage{orcidlink}
\usepackage[margin=0.75in]{geometry}
\usepackage{amsmath,amssymb,amsfonts,bm,mathtools}
\usepackage{booktabs,longtable,array,multirow}
\usepackage{xcolor}
\usepackage{hyperref}
\usepackage[T1]{fontenc}
\usepackage{lmodern}

\setlength{\parindent}{0pt}
\setlength{\parskip}{0.55em}
\renewcommand{\arraystretch}{1.22}

\newcommand{\PR}{\mathrm{PR}}
\newcommand{\eff}{\mathrm{eff}}
\newcommand{\Kerr}{\mathrm{Kerr}}
\newcommand{\GR}{\mathrm{GR}}
\newcommand{\obs}{\mathrm{obs}}
\newcommand{\src}{\mathrm{src}}
\newcommand{\vac}{\mathrm{vac}}
\newcommand{\hom}{\mathrm{hom}}
\newcommand{\TT}{\mathrm{TT}}
\newcommand{\Teuk}{\mathrm{Teuk}}
\newcommand{\dd}{\mathrm{d}}
\newcommand{\Tr}{\mathrm{Tr}}
\newcommand{\diag}{\mathrm{diag}}

\title{Projection Relativity:\\Complete Master Equation Sheet}
\author{Michael Stanislaus Oshetski \small \orcidlink{0009-0007-3623-7586} \\ 
\small Independent Researcher \\ 
\small ORCID: 0009-0007-3623-7586}\\
\date{May 2026}

\begin{document}
\maketitle

\section*{Purpose and Scope}

This master equation sheet consolidates the current Projection Relativity (PR) framework into one detailed mathematical reference aligned with the main manuscript. Its purpose is to provide a complete reconstruction map for the theory: the variables, operators, projection sectors, recovery conditions, spectral constraints, geometric closure relations, and observational signatures are collected in one place.

The organizing principle is that observable physical structure emerges through constrained projections of a deeper internal spectral manifold. Gravity, Higgs structure, electromagnetism, strong-field continuum corrections, and cosmological response are treated as observable sectors of the same projection field.

\begin{equation}
\boxed{
\Psi
\rightarrow
O_X
\rightarrow
\{u_n,\lambda_n\}
\rightarrow
P_i[\Psi]
\rightarrow
O_i
}
\end{equation}

This sheet uses the following status tags:
\[
D=\text{derived from the framework},\qquad
N=\text{normalization or convention},\qquad
C=\text{constrained by recovery/stability},\qquad
O=\text{future extension}.
\]


\section{Complete Variable and Parameter Dictionary}

\begin{longtable}{>{\raggedright\arraybackslash}p{0.16\textwidth}
>{\raggedright\arraybackslash}p{0.24\textwidth}
>{\raggedright\arraybackslash}p{0.47\textwidth}
>{\centering\arraybackslash}p{0.06\textwidth}}
\toprule
\textbf{Symbol} & \textbf{Name} & \textbf{Meaning in PR} & \textbf{Tag} \\
\midrule
\endhead

$\Psi(x,\xi)$ & Master projection field & Underlying spectral field depending on observable spacetime coordinate $x$ and internal coordinate $\xi$. All observable sectors are projections of this field. & D \\

$x^\mu$ & Spacetime coordinate & Effective low-energy coordinates in which projected observables are measured. & N \\

$\xi,w$ & Internal coordinates & Coordinates on the internal spectral manifold. In the reduced one-dimensional validation layer, $w$ is used. & D/N \\

$\mathcal H_P$ & Projection Hilbert space & Internal Hilbert space in which $\Psi$, $O_X$, eigenmodes, and projection operators are defined. & D \\

$O_X$ & Internal spectral operator & Hamiltonian-like internal operator generating the spectral eigenbasis and gap structure. & D \\

$V_X(w)$ & Internal potential & Confining projection potential defining bounded spectral structure and finite-core recovery scale. & C \\

$u_n(w)$ & Eigenmodes & Internal spectral modes satisfying $O_Xu_n=\lambda_nu_n$. & D \\

$\lambda_n$ & Eigenvalues & Internal spectral values associated with $u_n$. & D \\

$\mu_n^2$ & Excitation scales & Spectral gaps relative to the ground state; often $\mu_n^2=\lambda_n-\lambda_0$. & D \\

$\mu_{\min}^2$ & First spectral gap & Minimum internal excitation scale $\lambda_1-\lambda_0$; controls finite-core and propagator suppression. & D/C \\

$P_i[\Psi]$ & Projection operators & Sector-specific maps from the master field to observable structures. & D \\

$O_i$ & Observables & Projected physical observables such as $g_{\mu\nu}$, $H$, $A_\mu$, $\Sigma_X$, and $H_{\eff}$. & D \\

$P_g$ & Gravity projection & Stiffness projection producing metric structure and curvature response. & D \\

$P_H$ & Higgs projection & Displacement projection producing Higgs-sector scalar structure. & D \\

$P_A$ & Electromagnetic projection & Compact phase projection producing gauge-like electromagnetic structure. & D \\

$P_X$ & Strong-field correction projection & Continuum/self-energy projection producing weak strong-field residuals. & D/C \\

$P_\theta$ & Cosmological phase projection & Homogeneous low-energy phase projection associated with expansion response. & D/C \\

$Q_{\mu\nu}$ & Metric seed tensor & Bilinear projection correlation used to construct the effective spacetime metric. & D \\

$g_{\mu\nu}^{\eff}$ & Effective metric & Determinant-normalized observable metric recovered from $Q_{\mu\nu}$. & D \\

$g_{\mu\nu}^{\PR}$ & PR metric & Projection Relativity metric; equals Kerr outside $r_c$ in rotating exterior recovery. & D/C \\

$g_{\mu\nu}^{\Kerr}$ & Kerr metric & Required exterior rotating recovery target. & C \\

$G_P$ & Projection Newton coupling & Low-energy gravitational coupling; constrained by $G_P=G_N$. & C \\

$R_{\max}$ & Maximum curvature scale & Curvature saturation scale linked to spectral cutoff. & D/C \\

$r_c$ & Finite-core radius & Radius of saturated projection core replacing divergent singularity. & D/C \\

$\rho_P(\mu^2)$ & Spectral density & Weighted internal excitation density controlling propagator self-energy. & D \\

$F(k^2)$ & Projection self-energy & Spectral correction kernel in the propagator. & D/C \\

$D_{\mu\nu\rho\sigma}^{(P)}$ & Projection propagator & Effective spin-2 propagator corrected by $F(k^2)$. & D/C \\

$\Sigma_X$ & Strong-field correction & Weak projection self-energy correction to Kerr/Teukolsky response. & D/C \\

$A_X/A_0$ & Sideband amplitude ratio & Relative strength of projection-sector ringdown residual. & C \\

$A(x)$ & Projection amplitude & Amplitude/displacement coordinate in $\Psi=Ae^{i\theta}$; associated with Higgs sector. & D \\

$\theta(x)$ & Compact phase & Internal compact phase coordinate; associated with electromagnetic sector. & D \\

$A_\mu$ & Gauge potential & Electromagnetic phase projection. & D \\

$D_\mu$ & Gauge-covariant derivative & Covariant derivative $\partial_\mu-iqA_\mu$ generated by compact phase structure. & D \\

$F_{\mu\nu}$ & EM field strength & Gauge field tensor $\partial_\mu A_\nu-\partial_\nu A_\mu$. & D \\

$J_\theta^\mu$ & Phase current & Conserved Noether current of compact phase sector. & D \\

$Z_A$ & Compact closure factor & Geometric normalization yielding $\alpha_{\PR}^{-1}=4\pi Z_A$. & D/C \\

$\alpha_{\PR}$ & Geometric fine-structure constant & Emergent compact phase coupling scale. & D/C \\

$Z_H$ & Higgs normalization & Spectral overlap normalization for Higgs displacement sector. & D/C \\

$H(x)$ & Higgs projection observable & Scalar displacement projection $P_H[\Psi]$. & D \\

$H_{\PR}$ & PR Hamiltonian density & Unified projected energy density. & D \\

$S_{\PR}$ & Master action & Variational action over spacetime and internal coordinates. & D \\

$\theta_{\hom}$ & Homogeneous phase mode & Large-scale low-energy mode associated with cosmological expansion. & D/C \\

$H_{\eff}$ & Effective Hubble rate & Expansion response generated by homogeneous phase drift. & D/C \\

$B_{\PR}$ & Residual phase coherence scale & Predicted nanogauss-scale compact phase coherence observable. & C \\

\bottomrule
\end{longtable}


\section{Master Chain and Projection Correspondence}

The universal structure of PR is the projection chain
\begin{equation}
\boxed{
\Psi
\longrightarrow
O_X
\longrightarrow
\{u_n,\lambda_n\}
\longrightarrow
P_i[\Psi]
\longrightarrow
O_i
}
\end{equation}

The full observable-sector correspondence is
\begin{equation}
\boxed{
\Psi
\rightarrow
O_X
\rightarrow
\{u_n,\lambda_n\}
\rightarrow
\{P_g,P_H,P_A,P_X,P_\theta\}
\rightarrow
\{g_{\mu\nu},H,A_\mu,\Sigma_X,H_{\eff}\}.
}
\end{equation}

Equivalently,
\begin{equation}
\boxed{
O_{\phys}
=
\mathcal P\!\left[
\Psi,O_X,\{u_n,\lambda_n\}
\right].
}
\end{equation}

The native interpretation is:
\[
\text{gravity}=\text{stiffness projection},\qquad
\text{Higgs}=\text{displacement projection},\qquad
\text{electromagnetism}=\text{compact phase projection}.
\]


\section{Internal Hilbert Space}

The master field belongs to an internal projection Hilbert space:
\begin{equation}
\Psi\in\mathcal H_P.
\end{equation}

The internal inner product is
\begin{equation}
\langle f|g\rangle
=
\int d\xi\, f^*(\xi)g(\xi).
\end{equation}

For the reduced one-dimensional internal coordinate:
\begin{equation}
\mathcal H_w=L^2(\mathbb R,dw),
\qquad
\langle f|g\rangle
=
\int_{\mathbb R}dw\,f^*(w)g(w).
\end{equation}

The spectral basis satisfies
\begin{equation}
O_X|u_n\rangle=\lambda_n|u_n\rangle,
\end{equation}
\begin{equation}
\langle u_m|u_n\rangle=\delta_{mn},
\end{equation}
and
\begin{equation}
\boxed{
\sum_n|u_n\rangle\langle u_n|=\mathbf 1.
}
\end{equation}

The projection field expands as
\begin{equation}
\boxed{
\Psi(x,\xi)
=
\sum_n c_n(x)u_n(\xi).
}
\end{equation}


\section{Master Spectral Operator}

The internal spectral operator is
\begin{equation}
\boxed{
O_X
=
-\frac{d^2}{dw^2}
+
V_X(w).
}
\end{equation}

The validation-layer potential is
\begin{equation}
V_X(w)
=
\Lambda_X^2
\left(
1+a_2w^2+a_4w^4
\right),
\qquad
a_2>0,\quad a_4>0.
\end{equation}

Its curvature is
\begin{equation}
\frac{d^2V_X}{dw^2}
=
\Lambda_X^2
\left(
2a_2+12a_4w^2
\right)>0.
\end{equation}

The eigenvalue equation is
\begin{equation}
O_Xu_n(w)=\lambda_nu_n(w).
\end{equation}

The spectral gaps are
\begin{equation}
\Delta\lambda_n=\lambda_{n+1}-\lambda_n.
\end{equation}

The first spectral gap is
\begin{equation}
\boxed{
\mu_{\min}^2=\lambda_1-\lambda_0.
}
\end{equation}

The reference validation values are
\begin{equation}
\lambda_0\simeq2.322863529580,
\qquad
\lambda_1\simeq5.375830272676,
\end{equation}
giving
\begin{equation}
\mu_{\min}^2\simeq3.052966743096.
\end{equation}

The spectral gap propagates into the recovery hierarchy:
\begin{equation}
\boxed{
\mu_{\min}^2
\rightarrow
R_{\max}
\rightarrow
r_c
\rightarrow
F(k^2).
}
\end{equation}


\section{Master Action}

The full PR framework is generated by the spectral--projection action
\begin{equation}
\boxed{
S_{\PR}
=
\int d^4x\,d\xi\,
\sqrt{-g_{\eff}}
\left[
\frac12G^{AB}D_A\Psi^\dagger D_B\Psi
-
V_{\PR}(\Psi,\xi)
+
\mathcal L_{\src}
\right].
}
\end{equation}

Here:
\[
G^{AB}=\text{internal manifold metric},\qquad
D_A=\text{projection covariant derivative},
\]
\[
V_{\PR}=\text{internal projection potential},\qquad
\mathcal L_{\src}=\text{observable source sector}.
\]

The Euler--Lagrange equation schematically takes the form
\begin{equation}
\frac{1}{\sqrt{-g_{\eff}}}
D_A
\left(
\sqrt{-g_{\eff}}\,G^{AB}D_B\Psi
\right)
+
\frac{\partial V_{\PR}}{\partial\Psi^\dagger}
=
\frac{\partial\mathcal L_{\src}}{\partial\Psi^\dagger}.
\end{equation}

Thus the spectral structure, projection dynamics, source coupling, propagator response, and emergent metric are all treated as consequences of the same action principle.


\section{Emergent Metric and Weak Gravity Recovery}

The metric seed tensor is
\begin{equation}
Q_{\mu\nu}
=
\left\langle
\partial_\mu\Psi
\partial_\nu\Psi
\right\rangle.
\end{equation}

A more explicit internal metric form is
\begin{equation}
Q_{\mu\nu}
=
\left\langle
\partial_\mu\Psi^A
\partial_\nu\Psi^B
\right\rangle
G_{AB}.
\end{equation}

The observable metric is determinant normalized:
\begin{equation}
\boxed{
g_{\mu\nu}^{\eff}
=
\frac{Q_{\mu\nu}}{|\det Q|^{1/4}}.
}
\end{equation}

The weak-field expansion is
\begin{equation}
g_{\mu\nu}^{\eff}
=
\eta_{\mu\nu}+h_{\mu\nu},
\qquad
|h_{\mu\nu}|\ll1.
\end{equation}

The low-energy gravitational recovery condition is
\begin{equation}
G_{\mu\nu}^{\PR}
\rightarrow
8\pi G_PT_{\mu\nu},
\end{equation}
with
\begin{equation}
\boxed{
G_P=G_N.
}
\end{equation}

The massless graviton and weak-field recovery constraints are
\begin{equation}
\boxed{
F(0)=0,
\qquad
\left|
\frac{F(k^2)}{k^2}
\right|
\ll1
\quad
(k^2\ll\mu_{\min}^2).
}
\end{equation}

% ============================================================
\section{Finite-Core Regularization}

Classical singular behavior is replaced by bounded spectral saturation. The finite-core radius is
\begin{equation}
\boxed{
r_c
=
\left(
\frac{G_PM}{c^2R_{\max}}
\right)^{1/3}.
}
\end{equation}

The core hierarchy is
\begin{equation}
\mu_{\min}^2
\rightarrow
R_{\max}
\rightarrow
r_c.
\end{equation}

The interpretation is:
\[
r<r_c:\ \text{saturated projection core},
\qquad
r>r_c:\ \text{GR/Kerr exterior recovery}.
\]

The curvature is bounded schematically by
\begin{equation}
|R|\lesssim R_{\max}.
\end{equation}


\section{Kerr Exterior Recovery}

The rotating exterior recovery condition is
\begin{equation}
\boxed{
g_{\mu\nu}^{\PR}
=
g_{\mu\nu}^{\Kerr},
\qquad
r>r_c.
}
\end{equation}

The Kerr structure functions are
\begin{equation}
\Sigma_{\Kerr}
=
r^2+a^2\cos^2\theta,
\end{equation}
\begin{equation}
\Delta_{\Kerr}
=
r^2-\frac{2G_PM}{c^2}r+a^2.
\end{equation}

The PR structure functions are
\begin{equation}
\Sigma_{\PR}
=
\Sigma_{\Kerr}+\delta\Sigma(r),
\end{equation}
\begin{equation}
\Delta_{\PR}
=
\Delta_{\Kerr}+\delta_{\core}(r).
\end{equation}

Exterior recovery requires
\begin{equation}
\delta\Sigma(r)\rightarrow0,
\qquad
\delta_{\core}(r)\rightarrow0,
\qquad
r\gg r_c.
\end{equation}


\section{Projection Propagator and Self-Energy}

The spectral density is
\begin{equation}
\boxed{
\rho_P(\mu^2)
=
\sum_n
\left|
\left\langle
u_n
\middle|
\hat V
\middle|
u_0
\right\rangle
\right|^2
\delta(\mu^2-\mu_n^2).
}
\end{equation}

Normalization:
\begin{equation}
\boxed{
\int d\mu^2\,\rho_P(\mu^2)=1.
}
\end{equation}

The self-energy kernel is
\begin{equation}
\boxed{
F(k^2)
=
\int d\mu^2
\frac{\rho_P(\mu^2)}
{k^2-\mu^2+i\epsilon}.
}
\end{equation}

The projection propagator is
\begin{equation}
\boxed{
D_{\mu\nu\rho\sigma}^{(P)}(k)
=
\frac{
P_{\mu\nu\rho\sigma}^{(2)}
}{
k^2-F(k^2)+i\epsilon
}.
}
\end{equation}

The stability requirement is
\begin{equation}
k^2-F(k^2)>0
\end{equation}
in the physical low-energy recovery regime.


\section{Gravitational Waves and Ringdown}

Projection gravitational waves are perturbations of the emergent metric:
\begin{equation}
g_{\mu\nu}^{\eff}
=
\eta_{\mu\nu}
+
h_{\mu\nu}.
\end{equation}

The transverse-traceless propagation equation is
\begin{equation}
\boxed{
\left[
\Box+\Pi_X(\Box)
\right]
h_{\mu\nu}^{\TT}
=
16\pi G_PT_{\mu\nu}^{\TT}.
}
\end{equation}

The Kerr--Teukolsky inheritance structure is
\begin{equation}
\boxed{
O_{\PR}
=
O_{\Teuk}
+
\Sigma_X.
}
\end{equation}

The weak correction condition is
\begin{equation}
|\Sigma_X|\ll|O_{\Teuk}|.
\end{equation}

The ringdown response may be decomposed as
\begin{equation}
A(\omega)
=
A_{\Kerr}(\omega)+A_X(\omega),
\end{equation}
with
\begin{equation}
A_X(\omega)
=
\int d\mu\,
\frac{\rho_X(\mu)}
{\omega-\mu+i\Gamma_X(\mu)}.
\end{equation}

The projected sideband amplitude scale is
\begin{equation}
\boxed{
A_X/A_0\sim10^{-3}.
}
\end{equation}


\section{Higgs Displacement Sector}

The amplitude--phase decomposition is
\begin{equation}
\boxed{
\Psi(x)=A(x)e^{i\theta(x)}.
}
\end{equation}

Differentiation gives
\begin{equation}
\partial_\mu\Psi
=
e^{i\theta}
\left(
\partial_\mu A+iA\partial_\mu\theta
\right).
\end{equation}

The kinetic norm separates:
\begin{equation}
\boxed{
|\partial_\mu\Psi|^2
=
(\partial_\mu A)(\partial^\mu A)
+
A^2(\partial_\mu\theta)(\partial^\mu\theta).
}
\end{equation}

The Higgs projection is
\begin{equation}
\boxed{
H(x)=P_H[\Psi]=A(x).
}
\end{equation}

The Higgs-sector Lagrangian is
\begin{equation}
\mathcal L_H
=
\frac12(\partial_\mu A)(\partial^\mu A)
-
V(A).
\end{equation}

The potential is
\begin{equation}
V(A)=\alpha_HA^2+\beta_HA^4,
\qquad
\beta_H>0.
\end{equation}

The vacuum condition is
\begin{equation}
\frac{dV}{dA}
=
2\alpha_HA+4\beta_HA^3=0,
\end{equation}
giving
\begin{equation}
A_0^2
=
-\frac{\alpha_H}{2\beta_H}.
\end{equation}

The Higgs spectral normalization is
\begin{equation}
\boxed{
Z_H
=
\int d\xi\,u_0(\xi)P_H[\Psi].
}
\end{equation}

The Higgs mass curvature relation is
\begin{equation}
\boxed{
m_H^2
=
\left.
\frac{d^2V}{dA^2}
\right|_{A=A_0}
=
2\alpha_H+12\beta_HA_0^2.
}
\end{equation}


\section{Electromagnetic Compact Phase Sector}

The electromagnetic projection is
\begin{equation}
\boxed{
A_\mu=P_A[\Psi]\sim\partial_\mu\theta.
}
\end{equation}

The compact phase identification is
\begin{equation}
\boxed{
\theta\sim\theta+2\pi n,
\qquad
n\in\mathbb Z.
}
\end{equation}

The gauge-covariant derivative is
\begin{equation}
\boxed{
D_\mu=\partial_\mu-iqA_\mu.
}
\end{equation}

The compact gauge transformation is
\begin{equation}
\boxed{
A_\mu\rightarrow A_\mu+q^{-1}\partial_\mu\alpha.
}
\end{equation}

The electromagnetic field tensor is
\begin{equation}
\boxed{
F_{\mu\nu}
=
\partial_\mu A_\nu-\partial_\nu A_\mu.
}
\end{equation}

The Maxwell-type Lagrangian is
\begin{equation}
\mathcal L_{\EM}
=
-\frac14F_{\mu\nu}F^{\mu\nu}
+
J_\mu A^\mu.
\end{equation}

The Maxwell recovery equation is
\begin{equation}
\boxed{
\partial_\mu F^{\mu\nu}=J^\nu.
}
\end{equation}

The compact phase Lagrangian is
\begin{equation}
\mathcal L_A
=
\frac12A^2(\partial_\mu\theta)(\partial^\mu\theta).
\end{equation}

The Noether current is
\begin{equation}
\boxed{
J_\theta^\mu=A^2\partial^\mu\theta.
}
\end{equation}

The conservation law is
\begin{equation}
\boxed{
\partial_\mu J_\theta^\mu=0.
}
\end{equation}


\section{Geometric Fine-Structure Closure}

The fine-structure closure chain is
\begin{equation}
\boxed{
O_X
\rightarrow
\{u_n,\lambda_n\}
\rightarrow
\mu_{\min}
\rightarrow
P_A[\Psi]
\rightarrow
Z_A
\rightarrow
\alpha_{\PR}^{-1}.
}
\end{equation}

The compact closure factor is
\begin{equation}
Z_A
=
\oint d\theta\,u_0(\theta)P_A[\Psi].
\end{equation}

The geometric closure relation is
\begin{equation}
\boxed{
\alpha_{\PR}^{-1}=4\pi Z_A.
}
\end{equation}

The validated effective value is
\begin{equation}
\boxed{
\alpha_{\PR}^{-1}\simeq137.0365.
}
\end{equation}

The stability chain is
\begin{equation}
\delta\alpha
\rightarrow
\delta F_{\eff}(k^2)
\rightarrow
\delta\Sigma_X
\rightarrow
\text{loss of spectral regularity}.
\end{equation}


\section{Unified Projection Energy}

The projection-sector Lagrangian is
\begin{equation}
\mathcal L_{\PR}
=
\frac12(\partial_\mu A)(\partial^\mu A)
+
\frac12A^2(\partial_\mu\theta)(\partial^\mu\theta)
-
V(A)
-
\frac12\alpha_gP_g[\Psi]^2.
\end{equation}

Canonical momenta:
\begin{equation}
\Pi_A=\frac{\partial\mathcal L_{\PR}}{\partial\dot A}=\dot A,
\end{equation}
\begin{equation}
\Pi_\theta=\frac{\partial\mathcal L_{\PR}}{\partial\dot\theta}=A^2\dot\theta.
\end{equation}

The Hamiltonian density is
\begin{equation}
H_{\PR}
=
\Pi_A\dot A+\Pi_\theta\dot\theta-\mathcal L_{\PR}.
\end{equation}

Explicitly:
\begin{equation}
\boxed{
H_{\PR}
=
\frac12\dot A^2
+
\frac12(\nabla A)^2
+
\frac12A^2\dot\theta^2
+
\frac12A^2(\nabla\theta)^2
+
V(A)
+
\frac12\alpha_gP_g[\Psi]^2.
}
\end{equation}

The sector decomposition is
\begin{equation}
\boxed{
H_{\PR}=H_g+H_H+H_{\EM}+H_{\vac}.
}
\end{equation}

The observable relativistic energy relation is
\begin{equation}
\boxed{
Mc^2
=
E_{\PR}
=
\int d^3x\,H_{\PR}.
}
\end{equation}


\section{Projection Unitarity and Information Preservation}

The internal spectral evolution equation is
\begin{equation}
i\frac{\partial\Psi}{\partial t}
=
\hat H_P\Psi.
\end{equation}

The projection Hamiltonian is
\begin{equation}
\hat H_P=O_X+\hat V_{\int}.
\end{equation}

The unitary evolution operator is
\begin{equation}
U(t)=e^{-i\hat H_Pt}.
\end{equation}

Unitarity:
\begin{equation}
\boxed{
U^\dagger U=I.
}
\end{equation}

Norm conservation:
\begin{equation}
\boxed{
\langle\Psi|\Psi\rangle=1.
}
\end{equation}

Collapse is interpreted as
\begin{equation}
\boxed{
\text{collapse}
\rightarrow
\text{spectral excitation redistribution}
\rightarrow
\text{finite-core saturation}.
}
\end{equation}


\section{Cosmological Projection Structure}

The cosmological chain is
\begin{equation}
\boxed{
\Psi
\rightarrow
\theta_{\hom}
\rightarrow
J_{\vac}
\rightarrow
H_{\eff}(t).
}
\end{equation}

The homogeneous phase Lagrangian is
\begin{equation}
\mathcal L_{\hom}
=
\frac12f_\theta^2\dot\theta_{\hom}^2
-
V(\theta_{\hom}).
\end{equation}

The homogeneous phase potential is
\begin{equation}
V(\theta_{\hom})
=
\Lambda_\theta^4
\left[
1-\cos\left(
\frac{\theta_{\hom}}{f_\theta}
\right)
\right].
\end{equation}

Small-displacement expansion:
\begin{equation}
V(\theta_{\hom})
\simeq
\frac12m_\theta^2\theta_{\hom}^2,
\qquad
m_\theta^2=\frac{\Lambda_\theta^4}{f_\theta^2}.
\end{equation}

The vacuum source rule is
\begin{equation}
J_{\vac}\propto\langle T^{00}\rangle_{\hom}.
\end{equation}

The source selection hierarchy is
\begin{equation}
\boxed{
\langle T^{00}\rangle_{\hom}\rightarrow\theta_{\hom},
\qquad
T^{00}-\langle T^{00}\rangle_{\hom}\rightarrow P_g.
}
\end{equation}

The emergent Hubble rate is
\begin{equation}
\boxed{
H_{\eff}=A_\theta\dot\theta_{\hom}.
}
\end{equation}

The cosmological response equation is
\begin{equation}
\boxed{
\ddot\theta_{\hom}
+
3H_{\eff}\dot\theta_{\hom}
+
\frac{dV}{d\theta_{\hom}}
=
J_{\vac}.
}
\end{equation}

The late-time recovery structure is
\begin{equation}
H_{\eff}(z)=H_{\GR}(z)+\delta H_{\PR}(z),
\qquad
\left|
\frac{\delta H_{\PR}}{H_{\GR}}
\right|\ll1.
\end{equation}


\section{Observational Structure}

The observational hierarchy is
\begin{equation}
\boxed{
\mu_{\min}
\rightarrow
F(k^2)
\rightarrow
\Sigma_X
\rightarrow
O_{\obs}.
}
\end{equation}

Weak-field recovery:
\begin{equation}
g_{\mu\nu}^{\PR}=g_{\mu\nu}^{\GR},
\qquad
r\gg r_c.
\end{equation}

Kerr recovery:
\begin{equation}
g_{\mu\nu}^{\PR}=g_{\mu\nu}^{\Kerr},
\qquad
r>r_c.
\end{equation}

Strong-field condition:
\begin{equation}
k^2\sim\mu_{\min}^2.
\end{equation}

Projection correction amplitude:
\begin{equation}
\boxed{
A_X/A_0\sim10^{-3}.
}
\end{equation}

Residual phase-coherence scale:
\begin{equation}
\boxed{
B_{\PR}\sim10^{-12}\,\mathrm T\sim10\,\mathrm{nG}.
}
\end{equation}


\section{Validation and Completion Status}

\begin{longtable}{>{\raggedright\arraybackslash}p{0.34\textwidth}
>{\raggedright\arraybackslash}p{0.28\textwidth}
>{\raggedright\arraybackslash}p{0.28\textwidth}}
\toprule
\textbf{Framework Sector} & \textbf{Status} & \textbf{Validation Logic} \\
\midrule
\endhead

Internal spectral operator & Closed at framework level & Eigenvalue problem and spectral gap define internal scale. \\

Emergent metric & Closed at framework level & Determinant-normalized metric seed recovers weak-field geometry. \\

Weak-field GR recovery & Constrained and validated & Requires $G_P=G_N$, $F(0)=0$, and $|F/k^2|\ll1$. \\

Kerr exterior recovery & Constrained and validated & Requires $g_{\mu\nu}^{\PR}=g_{\mu\nu}^{\Kerr}$ for $r>r_c$. \\

Finite-core structure & Derived/constrained & Spectral gap sets $R_{\max}$ and finite-core radius $r_c$. \\

Projection propagator & Closed at effective level & Spectral self-energy remains suppressed at low energy. \\

Higgs displacement sector & First-order closed & Amplitude projection generates displacement sector and mass curvature. \\

Electromagnetic phase sector & First-order closed & Compact phase projection generates gauge-like sector and Noether current. \\

Fine-structure closure & First-order geometric closure & $\alpha_{\PR}^{-1}=4\pi Z_A$ gives validated closure value. \\

Unitarity/information & Effective-level closed & Hilbert evolution preserves norm and redistributes information spectrally. \\

Cosmological projection & Effective regime & Homogeneous phase drift provides late-time projection response. \\

Strong-field sidebands & Predictive & Weak continuum correction scale $A_X/A_0\sim10^{-3}$. \\

Full matter-sector extension & Future development & Additional projection sectors may extend the framework. \\

Nonperturbative UV completion & Future development & Full high-energy completion remains an extension beyond the current effective architecture. \\

\bottomrule
\end{longtable}


\section{Single-Line Native Definitions}

\begin{longtable}{>{\raggedright\arraybackslash}p{0.25\textwidth}
>{\raggedright\arraybackslash}p{0.65\textwidth}}
\toprule
\textbf{Concept} & \textbf{PR Native Definition} \\
\midrule
\endhead

Spacetime & The low-energy metric projection of the internal spectral manifold. \\
Gravity & The stiffness projection of the internal manifold. \\
Higgs structure & The displacement/amplitude projection of the manifold. \\
Electromagnetism & The compact phase projection of the manifold. \\
Fine-structure constant & The compact geometric closure scale of the electromagnetic phase sector. \\
Black-hole interior & A finite saturated projection-core region rather than a divergent singularity. \\
Kerr exterior & The recovered rotating low-energy exterior limit outside $r_c$. \\
Ringdown correction & A weak internal continuum response added to the Kerr--Teukolsky structure. \\
Cosmological expansion & Homogeneous phase drift of the projection manifold. \\
Energy & The low-energy exterior projection of conserved internal spectral energy. \\

\bottomrule
\end{longtable}

\end{document}
